\documentclass[10pt,a4paper]{amsart}
\usepackage[margin=19mm]{geometry}
\usepackage{amsmath,amssymb,mathtools}
\usepackage[scr=rsfs,cal=euler]{mathalfa}
\usepackage{xfrac}
\usepackage[unicode,hidelinks,bookmarks=false]{hyperref}
\newcommand{\ie}{i.e.\ }
\newcommand{\cf}{cf.\ }
\newcommand{\resp}{respectively\ }
\newcommand{\Subsection}{\subsection}
\providecommand{\nobreakdash}{\nobreak\hbox{-}}
\setlength{\emergencystretch}{2em}
\multlinegap0pt
\usepackage{amssymb}
\newtheorem{lemma}{Lemma}
\def\EE{\mathbb{E}}
\def\RR{\mathbb{R}}
\def\cF{\mathcal{F}}
\newcommand{\lc}{\left(}
\newcommand{\lk}{\left[}
\newcommand{\lt}{\left }
\newcommand{\rc}{\right)}
\newcommand{\rk}{\right]}
\newcommand{\rt}{\right}
\title{Growth rates for the H\"older coefficients of the linear stochastic fractional heat equation with rough dependence in space}
\author{Chang Liu, Bin Qian, Ran Wang}
\date{Source: arXiv:2507.22379v3 (CC BY 4.0)}
\begin{document}
\maketitle

\noindent\emph{Source and licence.} Growth rates for the H\"older coefficients of the linear stochastic fractional heat equation with rough dependence in space. Version arXiv:2507.22379v3 (CC BY 4.0), distributed by arXiv under the Creative Commons Attribution 4.0 International licence, http://creativecommons.org/licenses/by/4.0/, as recorded in the arXiv licence metadata for this exact version and shown on the abstract page https://arxiv.org/abs/2507.22379v3. Full licence notice: \texttt{licenses/arxiv-2507.22379-CC-BY-4.0.txt}.\par\smallskip

\noindent\textit{Editorial scope.} Counted unit: the article's lemma lem:diff3 (source0.tex lines 815--828). The article's complete original proof of it is delivered with it (source0.tex lines 829--872, one \texttt{proof} environment of 44 lines). No mathematics is written, repaired or re-derived: the statement and the proof are the article's own lines, in the article's own notation, and no retained line is substituted or re-flowed.  Retained uncounted as the source-local context the proof needs: lines 490--492; lines 518--520; lines 539--548; lines 743--745; lines 780--787; lines 797--803.  Nothing was omitted from any retained span, so the package carries no omission record.  No retained line contains a citation, so the wrapper carries no bibliography and no bibliography entry.  No retained line contains a figure environment, an image-inclusion command or drawing code, so no image is delivered and the package declares no asset dependency.  Editorial formatting only, in addition: an AMS \texttt{amsart} preamble that restates the article's own declarations and macros used by the retained lines, namely the environments \texttt{lemma} on separate counters, together with the standard packages those lines need.  The retained environments print in this document as Lemma 1 in order of appearance, whereas the article prints the same items with the numbering of its own full document.  The complete native source of the article as distributed in the arXiv e-print for this exact version is bundled unchanged as \texttt{sources/182409/source0.tex}.
     \begin{equation}\label{eq H product}
    \langle\phi,\psi\rangle_{\mathcal H}:= \EE\big[ W(\phi)W(\psi)\big]=c_{H}\int_{\RR_{+}\times \RR} \cF \phi(s,\xi) \overline{\cF \psi(s,\xi)}  |\xi|^{1-2H}d\xi ds,
     \end{equation}

\begin{equation}\label{Fourier}
(\mathcal {F}G_{\alpha}(t,x-\cdot))(\xi)=e^{-t|\xi|^{\alpha}-i\xi x}.
\end{equation}

\begin{lemma}\label{metric} There exist some  positive constants $c_{2,1}$ and $c_{2, 2}$ such that for any  $(t,x)$, $(s,y)\in \RR_+\times\RR$,
\begin{equation}\label{equ-metric-d1}
c_{2,1}\tilde d_{1}((t,x),(s,y))\le  d_1 ((t,x), (s,y)) \le c_{2,2}\tilde d_{1}((t,x),(s,y)),
\end{equation}
  where
\begin{equation*}
\tilde d_{1}((t,x),(s,y))
:= |x-y|^{\frac{2H+\alpha-2}{2}}\wedge(t\wedge s)^{\frac{2H+\alpha-2}{2\alpha}}+|t-s|^{\frac{2H+\alpha-2}{2\alpha}}.
\end{equation*}
\end{lemma}

\begin{align}\label{eq elem}
1-\cos(x)\ge \frac{x^2}{4} \ \ \ \text{for } |x|\le \frac{\pi}{2},
\end{align}

\begin{equation}\label{L-h-dEu0}
\begin{split}
\int_{\frac{|x-y|\pi}{4h}}^{\frac{|x-y|\pi}{2h}} \bigl(1-\cos\xi\bigr)  \xi^{3-2H-\alpha} \,d\xi
&\geq \sum_{k=k_0}^{k_1} \int_{I_k} \bigl(1-\cos\xi\bigr) \xi^{3-2H-\alpha} \,d\xi \\
&\geq \frac{1}{2}\sum_{k=k_0}^{k_1} \int_{I_k} \xi^{3-2H-\alpha} \,d\xi \\
&\geq \frac{1}{4} \int_{\frac{(6k_0+1)\pi}{3}}^{\frac{(6k_1+5)\pi}{3}} \xi^{3-2H-\alpha} \,d\xi.
\end{split}
\end{equation}

\begin{equation}\label{L-h-dEu0-2}
\begin{split}
 \int_{\frac{(6k_0+1)\pi}{3}}^{\frac{(6k_1+5)\pi}{3}} \xi^{3-2H-\alpha}  d\xi 
= &\,  \frac{1}{4-2H-\alpha} \left[ \left(\frac{(6k_1+5)\pi}{3}\right)^{4-2H-\alpha} - \left(\frac{(6k_0+1)\pi}{3}\right)^{4-2H-\alpha} \right] \\
 \geq &\,  \frac{\pi^{4-2H-\alpha}}{4-2H-\alpha} \left[ \left(\frac{13}{32}\right)^{4-2H-\alpha} - \left(\frac{11}{32}\right)^{4-2H-\alpha} \right] \left(\frac{|x-y|}{h}\right)^{4-2H-\alpha}.
\end{split}
\end{equation}

 \begin{lemma}\label{lem:diff3} 
\begin{itemize}
\item[(a)] (Upper bound)   There exists a positive constant $c_{2,8}$ such that for all   $0 \leq \theta \leq \frac{2H + \alpha - 2}{\alpha}$,  $x, y \in \mathbb{R}$ and   $0 < \tau \leq t$,
 \begin{align}\label{U-h-dEu 3-0}
   d_{3, t,\tau}(x,y)\leq c_{2,8} \tau^{\frac{\theta}{2}} \left(|x-y|^{\frac{2H+\alpha-2-\alpha\theta}{2}}\wedge t^{\frac{2H+\alpha-2-\alpha\theta}{2\alpha}}\right).
 \end{align}
 \item[(b)] (Lower bound)
 There exists a positive constant $c_{2,9}$ such that for all $|x-y|\geq t^{\frac1{\alpha}}$
  and $0 < \tau \leq \left(\frac{3}{32}\right)^{\alpha}|x-y|^{\alpha}$,
\begin{equation}\label{U-h-dEu 3}
       d_{3, t,\tau}(x,y)\geq  c_{2,9} \tau^{\frac{2H+\alpha-2}{2\alpha}}.
  \end{equation}
  \end{itemize}
  \end{lemma}

 \begin{proof}
  The upper bound in  \eqref{U-h-dEu 3-0}   follows immediately from Lemma \ref{metric} and the Minkowski inequality. Now, we     prove  the lower bound  \eqref{U-h-dEu 3}.


Using the Fourier representation \eqref{eq H product} and \eqref{Fourier}, along with the independence of stochastic integrals over $[0,t]$ and $[t,t+\tau]$, we have
      \begin{equation}\label{eq-d3}
\begin{split}
d_{3, t,\tau}^2(x,y)
\ge&\, \mathbb{E} \Bigg[ \left| \int_{t}^{t+\tau}\int_{\mathbb{R}} \left[ G_{\alpha}(t+\tau - s,x - z) - G_{\alpha}(t+\tau - s,y - z) \right] W(ds,dz)\right|^2 \Bigg] \\
=&\,  c_{H}\int_t^{t+\tau}\int_{\RR}  \lt| 1-e^{i(x-y)\xi} \rt|^2   e^{-2(t+\tau-s)|\xi|^{\alpha}}    |\xi|^{1-2H} d\xi ds  \\
=&\, 2c_{H}\int_{\RR_+}\lk 1-e^{-2\tau\xi^{\alpha}}\rk\lk 1-\cos(\xi|x-y|)\rk \xi^{1-2H-\alpha}d\xi\\
=&\, 2c_{H}|x-y|^{2H+\alpha-2}\int_{\RR_+}\lk 1-\exp\lc-\frac{2\tau\tilde{\xi}^{\alpha}}{|x-y|^{\alpha}}\rc\rk
\lk 1-\cos(\tilde{\xi})\rk \tilde{\xi}^{1-2H-\alpha}d\tilde\xi\\
\ge &\, 2c_{H}|x-y|^{2H+\alpha-2} \int_{\frac{|x-y|\pi}{2\tau^{1/\alpha}}}^{\frac{|x-y|\pi}{\tau^{1/\alpha}}}\lk 1-\exp\lc-\frac{2\tau\xi^{\alpha}}{|x-y|^{\alpha}}\rc\rk
\lk 1-\cos(\xi)\rk \xi^{1-2H-\alpha}d\xi\\
\ge &\,   2  c_{H} \left(1- \exp\left(-2^{1-\alpha}\pi^{\alpha}\right)\right) |x-y|^{2H+\alpha-2} \int_{\frac{|x-y|\pi}{2\tau^{1/\alpha}}}^{\frac{|x-y|\pi}{\tau^{1/\alpha}}} \lk 1-\cos(\xi)\rk \xi^{1-2H-\alpha}d\xi,
 \end{split}
\end{equation}
 where the change of variables   $\xi \mapsto \tilde \xi/|x-y|$  is used in the  third-last step, and the elementary inequality
  \[
    1 - \exp\left(-\frac{2\tau\xi^\alpha}{|x-y|^\alpha}\right) \geq
    1- \exp\left(-2^{1-\alpha}\pi^{\alpha}\right),\ \ \text{if}\ \ \xi\geq \frac{|x-y|\pi}{2\tau^{1/\alpha}},
    \]
is used   in the last step.


For any  $  \xi \in  \left[\frac{|x-y|\pi}{2\tau^{1/\alpha}},  \frac{|x-y|\pi}{\tau^{1/\alpha}} \right]$, by \eqref{eq elem},  we have
    \[
  1 \ge    1 - \cos\left(\frac{\tau^{1/\alpha}\xi}{2|x-y|}\right) \ge \frac{\tau^{2/\alpha}\xi^2}{16|x-y|^2}.
    \]
    Consequently, when $\tau \leq \left(\frac{3}{32}\right)^\alpha   |x-y|^{\alpha}$, it holds that
     \begin{equation}\label{h-tau}
\begin{split}
   \int_{\frac{|x-y|\pi}{2\tau^{1/\alpha}}}^{\frac{|x-y|\pi}{\tau^{1/\alpha}}} \left( 1-\cos(\xi)\right) \xi^{1-2H-\alpha}d\xi
\ge &\,  \frac{1}{16}\tau^{\frac{2}{\alpha}}|x-y|^{-2}
\int_{\frac{|x-y|\pi}{2\tau^{1/\alpha}}}^{\frac{|x-y|\pi}{\tau^{1/\alpha}}}
\left( 1-\cos(\xi)\right) \xi^{3-2H-\alpha}d\xi\\
\geq&\, c_{2,10} \tau^{\frac{2H+\alpha-2}{\alpha}}|x-y|^{2-2H-\alpha},
      \end{split}
\end{equation}
   where $c_{2,10}>0$, and the last inequality follows from \eqref{L-h-dEu0} and \eqref{L-h-dEu0-2}.

Combining \eqref{eq-d3} and \eqref{h-tau}, we obtain \eqref{U-h-dEu 3}. The proof is complete.
 \end{proof}

\end{document}
